% TOP_PROBLEM: {"id": 30007068, "title": "Universal-cover products for integrable tangent summands", "category_id": 5, "category": {"id": 5, "name": "algebraic_geometry", "display_name": "Algebraic Geometry", "description": "Geometric objects defined by polynomial equations.", "slug": "algebraic-geometry", "order_index": 5, "created_at": "2026-07-31T15:26:25.671Z"}, "status": "open"}
% FIRST_POSTED: 2026-10-06
% ATTEMPT_STATUS: solved
% OPENAI_CLAIM: {"schema_version":1,"source_repo":"https://github.com/openai/math","source_commit":"adc7f1241b42e322a6451854ab7e4b4c146bf78a","release_date":"2026-10-06","match":"full","resolution":"proved","independently_reviewed":false,"second_pass":"agreed","families":[{"family":"052","title":"Tangent splittings and product decompositions","subject":"Algebraic and complex geometry","manuscripts":[{"title":"Universal-cover splitting for compact Kähler manifolds","pdf_path":"preprints/Universal-cover-splitting-for-compact-Kahler-manifolds-September-23-2026/paper.pdf","readme_path":"preprints/Universal-cover-splitting-for-compact-Kahler-manifolds-September-23-2026/README.md","date":"2026-09-23","theorem_ref":"preprints/Universal-cover-splitting-for-compact-Kahler-manifolds-September-23-2026/build/intro.tex#thm:main"}],"lean":{"doc_path":"lean/docs/052.md","comparators":["lean/ComparatorChallenges/KahlerSplitting.json","lean/ComparatorChallenges/KahlerSplitting.lean","lean/ComparatorChallenges/SplitTangentIntegrability.json","lean/ComparatorChallenges/SplitTangentIntegrability.lean"],"declarations":["OAI.UniversalCoverSplitting.main"],"covers_main_theorem":true},"reasoning_trace":null}]}
\documentclass[11pt]{article}
\usepackage[margin=25mm]{geometry}
\usepackage{fontspec}
\setmainfont{Latin Modern Roman}
\usepackage{amsmath,amssymb,mathtools}
\usepackage{unicode-math}
\setmathfont{Latin Modern Math}
\usepackage{enumitem}
\usepackage{xurl,hyperref,bookmark}
\hypersetup{colorlinks=true,urlcolor=blue}
\setcounter{secnumdepth}{0}
\setlength{\parindent}{0pt}
\setlength{\parskip}{5pt}
\setlength{\emergencystretch}{3em}
\begin{document}
\section*{Universal-cover products for integrable tangent summands}

\subsection{OpenAI's claim}
If a compact Kähler tangent bundle splits into two integrable positive-rank subbundles, its ordinary universal cover is claimed to split compatibly with those particular distributions. The cited manuscript argues for that result; its relation to this problem is specified below.

\subsection{Outline of the argument}
Manuscript 1 (family 052, 2026-09-23): Separate tangency identities remove fibers of the meromorphic graph when a box is extended across a Hartogs figure. Kähler area bounds on annuli then extend adapted boxes across convex boundaries; continuation produces the two global leaf factors with the original summands as their tangent bundles.

\subsection{Scope relative to this problem}
Manuscript 1 (family 052, 2026-09-23): The problem selects the mathematical statement summarized here, including its stated hypotheses and parameter range. The cited theorem gives a claimed proof of that assertion in full.

\subsection{Formal verification}
For family 052, the scope page documents the following selections. One selected statement splits the universal cover assuming integrability of both supplied tangent summands. Another establishes that integrability for rationally connected projective manifolds; its compatible-product consequence is not itself that selected statement. The documented statement covers the main claim of manuscript 1. Comparator files: KahlerSplitting.lean, SplitTangentIntegrability.lean. Declarations: OAI.UniversalCoverSplitting.main, OAI.SplitTangent.integrability\_main.

This site has not executed the Lean files or checked their dependencies. These descriptions report the scope of the linked declarations, rather than a new verification of the proofs.

\subsection{Status caveat}
OpenAI attributes these results to an internal model. Its release notice cautions that results lacking full formalization can contain mistakes. This site has not independently assessed the proofs; a solved label records the scope of the reported claim.

\subsection{Sources}
\begin{itemize}
\item OpenAI, \emph{Universal-cover splitting for compact Kähler manifolds}, 2026-09-23; citation key \texttt{OAI:Universal-cover-splitting-for-compact-Kahler-manifolds-September-23-2026}. \href{https://github.com/openai/math/blob/adc7f1241b42e322a6451854ab7e4b4c146bf78a/preprints/Universal-cover-splitting-for-compact-Kahler-manifolds-September-23-2026/paper.pdf}{PDF (GitHub)}; \href{https://github.com/openai/math/raw/adc7f1241b42e322a6451854ab7e4b4c146bf78a/preprints/Universal-cover-splitting-for-compact-Kahler-manifolds-September-23-2026/paper.pdf}{PDF (direct)}; \href{https://github.com/openai/math/blob/adc7f1241b42e322a6451854ab7e4b4c146bf78a/preprints/Universal-cover-splitting-for-compact-Kahler-manifolds-September-23-2026/README.md}{Manuscript page and citation}; \href{https://github.com/openai/math/blob/main/preprints/Universal-cover-splitting-for-compact-Kahler-manifolds-September-23-2026/paper.pdf}{Latest version}.
\item Family 052: \href{https://github.com/openai/math/blob/adc7f1241b42e322a6451854ab7e4b4c146bf78a/lean/docs/052.md}{Lean scope}.
\item \href{https://github.com/openai/math/blob/adc7f1241b42e322a6451854ab7e4b4c146bf78a/lean/ComparatorChallenges/KahlerSplitting.json}{Comparator KahlerSplitting.json}.
\item \href{https://github.com/openai/math/blob/adc7f1241b42e322a6451854ab7e4b4c146bf78a/lean/ComparatorChallenges/KahlerSplitting.lean}{Comparator KahlerSplitting.lean}.
\item \href{https://github.com/openai/math/blob/adc7f1241b42e322a6451854ab7e4b4c146bf78a/lean/ComparatorChallenges/SplitTangentIntegrability.json}{Comparator SplitTangentIntegrability.json}.
\item \href{https://github.com/openai/math/blob/adc7f1241b42e322a6451854ab7e4b4c146bf78a/lean/ComparatorChallenges/SplitTangentIntegrability.lean}{Comparator SplitTangentIntegrability.lean}.
\item \href{https://github.com/openai/math/blob/adc7f1241b42e322a6451854ab7e4b4c146bf78a/CONTENTS.md}{OpenAI catalogue entry}.
\end{itemize}
\end{document}
